% GENERATED FILE. Edit canonical lessons, then run npm run generate:books.
% canonical-source-hash: fnv1a64:26538b29f60ae2f6
% canonical-lessons: ML-C105-raajyam, ML-C105-nagaram, ML-C105-thalasthaanam, ML-R105-places-recall

\chapter{Places}
\label{ch:ml-places}
\hlchaptermodality{105}

\begin{chapteropening}
\textbf{By the end of this chapter:} \emph{I can name a country, a town, a city and a capital, and say what kind of place Kerala and India are.}

From cold: the scale of places, and the one word on it that does not sort by size.
\end{chapteropening}

\section[rājyaṁ]{\ml{രാജ്യം} — a country}

\label{lesson:ML-C105-raajyam}



\begin{quote}
Say \textbf{\ml{ഇന്ത്യ}}. Say \textbf{\ml{കേരളം}}. You have had both of those names for a long while — and no word for \textbf{what kind of thing} either one is.
\end{quote}

\subsection*{What to know first}
\textbf{\ml{രാജ്യം}} (\emph{rājyaṁ}) — \textbf{a country}.

\textbf{\ml{ഇന്ത്യ}} is one. \textbf{\ml{കേരളം}} is not — it is a part of one. You have been able to say both names since long before you had a word for the difference between them.

\begin{culture}
The Sanskrit behind it means \textbf{a realm}, and it is built on the word for a \textbf{king}. A country is named, in this word, for \textbf{who once ruled it} — not for its land, its people or its borders.

\textbf{It ends in -\ml{ം}}, which by now tells you what it will do when a case ending arrives.
\end{culture}

\subsection*{Your turn}
Say these aloud:

\begin{itemize}
\raggedright
  \item \emph{rājyaṁ}
  \item \emph{indya}, then \emph{rājyaṁ} — the name, then what kind of thing it is
  \item what the Sanskrit behind it is built on
\end{itemize}

\begin{tcolorbox}[breakable,colback=teal!4,colframe=teal!35!black,title={Before you move on}]
A country? (\textbf{\ml{രാജ്യം}}.) Which country have you been able to name since long before you had this word? (\textbf{\ml{ഇന്ത്യ}}.) What is the word built on? (\textbf{The word for a king}.)
\end{tcolorbox}
\section[nagaraṁ]{\ml{നഗരം} — a city}

\label{lesson:ML-C105-nagaram}



\begin{quote}
Say \emph{a country}. Then say the word you were given for a settled place somebody is from.
\end{quote}

\subsection*{What to know first}
\textbf{\ml{നഗരം}} (\emph{nagaraṁ}) — \textbf{a city}.

\noindent
\begin{tabularx}{\linewidth}{@{}>{\raggedright\arraybackslash}X>{\raggedright\arraybackslash}X>{\raggedright\arraybackslash}X@{}}
\toprule
\textbf{} & \textbf{} & \textbf{} \\
\midrule
\textbf{\ml{രാജ്യം}} & \emph{rājyaṁ} & a country \\
\textbf{\ml{നഗരം}} & \emph{nagaraṁ} & a city \\
\textbf{\ml{ഊര്}} & \emph{ūrŭ} & a settled place — \textbf{yours already} \\
\bottomrule
\end{tabularx}

\begin{culture}
\textbf{The line between the bottom two is not as sharp as the table makes it look}, and it is worth saying so rather than teaching you a rule that Malayalam does not keep.

\textbf{\ml{നഗരം}} is \textbf{Sanskrit} and the more formal of the two. \textbf{\ml{ഊര്}} is \textbf{Dravidian}, old, and the everyday one — the word you were told is buried in place-names the length of the south. A large place will usually be called \textbf{\ml{നഗരം}}; a smaller one \textbf{\ml{ഊര്}} — but the same place can take either, and which one a speaker reaches for has as much to do with \textbf{register} as with size.

\textbf{What is reliable is the register, not the size.} If you want the formal word, you want this one.
\end{culture}

\subsection*{Your turn}
Say these aloud:

\begin{itemize}
\raggedright
  \item \emph{nagaraṁ}
  \item the three in order — \emph{rājyaṁ}, \emph{nagaraṁ}, \emph{ūrŭ}
  \item which of the two smaller words is the formal one, and where it came from
\end{itemize}

\begin{tcolorbox}[breakable,colback=teal!4,colframe=teal!35!black,title={Before you move on}]
A city? (\textbf{\ml{നഗരം}}.) Which of your two words is Sanskrit, and which Dravidian? (\textbf{\ml{നഗരം}} Sanskrit, \textbf{\ml{ഊര്}} Dravidian.) Is the difference between them strictly about size? (\textbf{No} — it is as much about how formal you are being.)
\end{tcolorbox}
\section[talasthānaṁ]{\ml{തലസ്ഥാനം} — a capital}

\label{lesson:ML-C105-thalasthaanam}



\begin{quote}
Say \emph{a city}. Then say the word for the part of your body that sits on top.
\end{quote}

\subsection*{What to know first}
\textbf{\ml{തലസ്ഥാനം}} (\emph{talasthānaṁ}) — \textbf{a capital}.

\textbf{It opens with a word you already have.} \emph{tala} is the \textbf{head} you learned among the body parts, and \emph{sthānaṁ} is a \textbf{place}. A capital is a \textbf{head-place} — and the second half is given here in sound only, because this book has not taught it as a word of its own.

\begin{culture}
\textbf{English says exactly the same thing, and hides it.} \emph{Capital} comes from Latin \emph{caput} — \textbf{head}. Two languages, two roots, one idea: the city at the top of the body.

Malayalam keeps the head visible. You can hear \emph{tala} in the word every time you say it; an English speaker has to be told.
\end{culture}

\subsection*{Your turn}
Say these aloud:

\begin{itemize}
\raggedright
  \item \emph{talasthānaṁ}, slowly — \emph{tala … sthānaṁ}
  \item the body part hiding at the front of it
  \item \emph{nagaraṁ}, then \emph{talasthānaṁ} — a city, then the head one
\end{itemize}

\begin{tcolorbox}[breakable,colback=teal!4,colframe=teal!35!black,title={Before you move on}]
A capital? (\textbf{\ml{തലസ്ഥാനം}}.) What are its two halves? (\emph{tala}, the \textbf{head}, and \emph{sthānaṁ}, a \textbf{place}.) What does the English word hide that the Malayalam one shows? (\textbf{The same head} — Latin \emph{caput}.)
\end{tcolorbox}
\section[rājyaṁ enna ōrmma]{(country, city, capital) — from cold}

\label{lesson:ML-R105-places-recall}



\begin{quote}
Close the lessons before this one. Say \emph{country}, \emph{city} and \emph{capital}.
\end{quote}

\begin{grammarlens}[title={Grammar Lens: a scale, and one word that is not on it}]
\noindent
\begin{tabularx}{\linewidth}{@{}>{\raggedright\arraybackslash}X>{\raggedright\arraybackslash}X>{\raggedright\arraybackslash}X@{}}
\toprule
\textbf{} & \textbf{} & \textbf{} \\
\midrule
\textbf{\ml{രാജ്യം}} & \emph{rājyaṁ} & a country \\
\textbf{\ml{നഗരം}} & \emph{nagaraṁ} & a city \\
\textbf{\ml{ഊര്}} & \emph{ūrŭ} & a settled place — yours since long before this chapter \\
\textbf{\ml{തലസ്ഥാനം}} & \emph{talasthānaṁ} & a capital \\
\bottomrule
\end{tabularx}

\textbf{The last one is not a size.} A capital is a city picked out by what happens in it, and it can be large or small. The first three sort places by how much ground they cover; this one sorts by \textbf{what a place does}.
\end{grammarlens}

\begin{grammarlens}[title={What you've built}]
\textbf{You had the names before you had most of the nouns.} \textbf{\ml{കേരളം}} and \textbf{\ml{ഇന്ത്യ}} have been yours for a long while, and \textbf{\ml{ഊര്}} and \textbf{\ml{നാട്}} gave you the small end of the scale and a word that floats over the rest of it. What was missing was the top — a word that does \emph{not} float — and the large settlement, and the capital.

And one of the four was half-built already: \textbf{\ml{തലസ്ഥാനം}} opens with \emph{tala}, the \textbf{head} you learned among the body parts. The word was waiting on a piece you had.
\end{grammarlens}

\subsection*{Your turn}
Say these aloud:

\begin{itemize}
\raggedright
  \item \emph{rājyaṁ}, \emph{nagaraṁ}, \emph{ūrŭ}, \emph{talasthānaṁ}
  \item which one is not a size
  \item which of them is the old Dravidian one
  \item \emph{indya}, then what kind of thing it is
\end{itemize}

\begin{tcolorbox}[breakable,colback=teal!4,colframe=teal!35!black,title={Before you move on}]
Country, city, capital? (\textbf{\ml{രാജ്യം}, \ml{നഗരം}, \ml{തലസ്ഥാനം}}.) Which of them does not sort places by size? (\textbf{\ml{തലസ്ഥാനം}}.) What is hiding at the front of it? (\emph{tala}, the \textbf{head}.)
\end{tcolorbox}
